\documentclass[11pt,a4paper]{article}

% -------------------------------------------------
% Packages
% -------------------------------------------------
\usepackage[margin=2.5cm]{geometry}
\usepackage{amsmath}
\usepackage{amssymb}
\usepackage{mathtools}
\usepackage{booktabs}
\usepackage{longtable}
\usepackage{array}
\usepackage{enumitem}
\usepackage{microtype}
\usepackage{listings}
\usepackage{xcolor}
\usepackage[hidelinks]{hyperref}

% -------------------------------------------------
% Custom commands
% -------------------------------------------------
\newcommand{\R}{\mathbb{R}}
\newcommand{\N}{\mathbb{N}}
\newcommand{\Tstar}{T^{\ast}}
\newcommand{\Acal}{\mathcal{A}}
\newcommand{\Fcal}{\mathcal{F}}
\newcommand{\Pcal}{\mathcal{P}}
\newcommand{\maxSeg}{\operatorname{maxSeg}}

\lstset{
  basicstyle=\ttfamily\small,
  breaklines=true,
  columns=fullflexible,
  frame=single,
  showstringspaces=false
}

\title{Mixed-Integer Programming Formulation and Implementation of the T-ASR Solver}
\author{}
\date{}

\begin{document}

\maketitle

\tableofcontents
\newpage

\section{Overview}

This project implements an experimental mixed-integer programming approach to
 the T-Adaptive Segment Routing (T-ASR) problem of the
\href{https://roadef.org/challenge/2026/en/index.php}{ROADEF/EURO 2026 ``Keep the Flow!'' challenge}.
The implementation first computes ECMP split coefficients for every time period,
segment, and available network arc. It then constructs a sparse mixed-integer
linear model whose routing decisions select segment paths for every demand and
time period.

The intended objective is lexicographic minimization of the vector of arc
utilizations over all available arc--time pairs. The first level minimizes the
maximum utilization. Subsequent levels are handled by a persistent bottleneck
peeling procedure whenever individual bottlenecks can be proved unavoidable.
If that proof fails, the implementation switches to a residual ordered-$k$-sum
formulation over the still-unfixed loads.

The mathematical model is represented through a solver-independent sparse
linear-model interface. The same formulation can be solved with HiGHS or the
Gurobi C API. The implementation also includes sparse preprocessing, an
optional common-routing warm start, persistent native solver sessions, MIP
starts for rebuilt residual models, global time-limit handling, graceful
interruption, physical-load consistency checks, LP export, and JSON solution
extraction.

\section{ECMP split-coefficient preprocessing}

The split coefficients are computed before the MILP is built. For each time
period, the implementation copies the network arc metrics and sets the metric
of every intervened arc to infinity. For each destination node, reverse Dijkstra
then computes shortest-path distances in the time-dependent network.

For an arc $a=(u,v)$ and destination $r$, the arc is eligible for forwarding
toward $r$ when
\[
d_t(u,r)\approx w_a^{(t)}+d_t(v,r),
\]
where $d_t(\cdot,r)$ is the shortest-path distance and $w_a^{(t)}$ is the
arc metric at time $t$. For every reachable source--destination pair, one unit
of flow is initialized at the source. At each node, the incoming flow is split
equally among all eligible forwarding arcs, and the resulting arc flows are
stored as the coefficients $r_{i,j,a,t}$. Only positive coefficients are kept
in the sparse representation.

This preprocessing models equal-cost multipath (ECMP) forwarding and makes the
coefficients fixed parameters of the subsequent optimization model.

\section{Sets and notation}

Let
\[
G=(V,A)
\]
be the directed network. Let
\[
T=\{0,1,\ldots,h-1\},
\qquad
\Tstar=T\setminus\{0\},
\]
where $T$ is the set of time periods. For every $t\in T$, let $q(t)\subseteq A$
be the set of unavailable arcs and define
\[
A_t:=A\setminus q(t).
\]
Each demand $d\in D$ has a source node $s_d$ and a target node $t_d$.
The demand volume at time $t$ is denoted by $\nu_{d,t}$.

For a segment $(i,j)$ and time $t$, the implementation computes sparse ECMP
split coefficients
\[
r_{i,j,a,t}\geq 0,
\]
where $r_{i,j,a,t}$ is the fraction of the flow associated with segment $(i,j)$
that traverses arc $a$ at time $t$. Intervened arcs are assigned infinite
forwarding metric and consequently do not receive positive split coefficients.

Define the set of reachable segments at time $t$ as
\[
P_t:=\left\{(i,j)\in V\times V:\ i\neq j,\;
 r_{i,j,a,t}>0\text{ for at least one }a\in A_t\right\}.
\]
Only segments in $P_t$ are considered when creating routing variables.

The set of available arc--time pairs is
\[
\Acal:=\{(t,a):t\in T,\ a\in A_t\},
\qquad m:=|\Acal|.
\]
The $m$ corresponding utilization values are the components of the
lexicographic objective.

\section{Parameters and numerical tolerances}

\begin{longtable}{
    >{\raggedright\arraybackslash}p{0.23\textwidth}
    >{\raggedright\arraybackslash}p{0.67\textwidth}
}
\toprule
\textbf{Parameter} & \textbf{Description} \\
\midrule
\endfirsthead
\toprule
\textbf{Parameter} & \textbf{Description} \\
\midrule
\endhead

$c_a$ & Capacity of arc $a\in A$. \\[0.4em]
$r_{i,j,a,t}$ & ECMP split coefficient for segment $(i,j)$, arc $a$, and time $t$. \\[0.4em]
$\nu_{d,t}$ & Traffic volume of demand $d$ at time $t$. \\[0.4em]
$\kappa_t$ & Maximum total number of segment-selection changes between $t-1$ and $t$. \\[0.4em]
$\maxSeg$ & Maximum number of segments in one demand path. \\[0.4em]
$U$ & Optional maximum-load upper bound obtained from the common-routing warm-start model. \\[0.4em]
$L_k$ & $k$-th ordered utilization level, with $L_1\geq L_2\geq\cdots$. \\[0.4em]
$S_k$ & Cumulative ordered objective, $S_k=\sum_{j=1}^{k}L_j$. \\[0.4em]
$R_q$ & Cumulative ordered objective of order $q$ on the residual unfixed-load vector. \\[0.4em]
$\varepsilon_{\mathrm{lex}}$ & Numerical tolerance used for preserving/fixing lexicographic levels; the implementation uses $10^{-6}$. \\[0.4em]
$\tau$ & Bottleneck proof/early-rejection tolerance; the implementation uses $10^{-7}$. \\
\bottomrule
\end{longtable}

The code also uses smaller tolerances for diagnostics and segment fixing. The
load equations are checked against physical utilization recomputed directly
from the selected binary routing variables.

\section{Decision variables}

\begin{longtable}{
    >{\raggedright\arraybackslash}p{0.25\textwidth}
    >{\raggedright\arraybackslash}p{0.65\textwidth}
}
\toprule
\textbf{Variable} & \textbf{Description} \\
\midrule
\endfirsthead
\toprule
\textbf{Variable} & \textbf{Description} \\
\midrule
\endhead

$x_{i,j}^{d,t}\in\{0,1\}$ & One if segment $(i,j)$ is selected in the path of demand $d$ at time $t$. It is created only for reachable segments that survive global maximum-load preprocessing. \\[0.4em]
$\lambda_{a,t}\in\R_{\geq0}$ & Exact utilization of available arc $a$ at time $t$. Its load equation is an equality, not an inequality. \\[0.4em]
$y_{i,j}^{d,t}\in[0,1]$ & Auxiliary variable for the absolute difference between the segment-selection variables at times $t-1$ and $t$. It is created only when the segment exists in at least one adjacent period. \\[0.4em]
$z\in\R_{\geq0}$ & Maximum-load variable used by the first and all persistent peeled min--max iterations. The same column is reused; no new $z_{\mathrm{free}}$ variable is created for each peeled level. \\[0.4em]
$\alpha_i\in\R$ & Threshold variable for residual ordered-sum level $i\geq2$. \\[0.4em]
$\beta_p^i\in\R_{\geq0}$ & Excess variable for residual free-load position $p$ and residual level $i\geq2$. Fixed bottleneck loads are not included in these auxiliary variables. \\
\bottomrule
\end{longtable}

\section{Base formulation}

The following constraints are present in every full T-ASR model. Terms for
variables that were not created during preprocessing are omitted.

\subsection{Exact arc-utilization equations}

For every $(t,a)\in\Acal$, the utilization is defined by
\begin{equation}
 c_a\lambda_{a,t}
 -\sum_{d\in D}\sum_{(i,j)\in P_t}
 r_{i,j,a,t}\,\nu_{d,t}\,x_{i,j}^{d,t}=0.
\tag{1}
\end{equation}

Thus $\lambda_{a,t}$ is the physical utilization induced by the routing
variables. The equality is important for the later bottleneck tests: a load
variable cannot retain artificial slack above the utilization implied by $x$.

\subsection{Flow conservation}

For every node, demand, and time period,
\begin{equation}
\sum_{j:(i,j)\in P_t}x_{i,j}^{d,t}
-
\sum_{j:(j,i)\in P_t}x_{j,i}^{d,t}
=
\begin{cases}
1, & i=s_d,\\
-1, & i=t_d,\\
0, & \text{otherwise},
\end{cases}
\quad
\forall i\in V,\ d\in D,\ t\in T.
\tag{2}
\end{equation}

\subsection{Node-degree inequalities}

The formulation limits branching and strengthens the path model through
\begin{align}
\sum_{j:(i,j)\in P_t}x_{i,j}^{d,t}&\leq1,
&&\forall i\in V,\ d\in D,\ t\in T, \tag{3a}\\
\sum_{j:(j,i)\in P_t}x_{j,i}^{d,t}&\leq1,
&&\forall i\in V,\ d\in D,\ t\in T. \tag{3b}
\end{align}
The final solution extractor additionally rejects branching, cycles, and
selected segments that do not form a valid source-to-target path.

\subsection{Maximum number of segments}

For every demand and time period,
\begin{equation}
\sum_{(i,j)\in P_t}x_{i,j}^{d,t}
\leq \maxSeg,
\qquad \forall d\in D,\ t\in T.
\tag{4}
\end{equation}

\subsection{Reconfiguration budget}

For $t\in\Tstar$, the intended budget is
\begin{equation}
\sum_{d\in D}\sum_{\substack{i,j\in V\\i\neq j}}
\left|x_{i,j}^{d,t}-x_{i,j}^{d,t-1}\right|
\leq\kappa_t.
\tag{5}
\end{equation}
A segment variable missing in one of the two periods is interpreted as zero.
The implementation linearizes (5) with
\begin{align}
y_{i,j}^{d,t}&\geq x_{i,j}^{d,t}-x_{i,j}^{d,t-1}, \tag{5a}\\
y_{i,j}^{d,t}&\geq x_{i,j}^{d,t-1}-x_{i,j}^{d,t}, \tag{5b}\\
\sum_{d\in D}\sum_{\substack{i,j\in V\\i\neq j}}y_{i,j}^{d,t}
&\leq\kappa_t. \tag{5c}
\end{align}
Since the $x$ variables are binary and $0\leq y\leq1$, these constraints make
the budget count the required changes.

\subsection{Variable domains}

\begin{align}
x_{i,j}^{d,t}&\in\{0,1\}, \tag{6a}\\
\lambda_{a,t}&\in\R_{\geq0}, \tag{6b}\\
0\leq y_{i,j}^{d,t}&\leq1. \tag{6c}
\end{align}

\section{Sparse preprocessing and variable bounds}

\subsection{Reachability preprocessing}

The implementation does not create all potential $|D||T||V|(|V|-1)$ segment
variables. A segment $(i,j)$ is omitted at time $t$ if its sparse split-coefficient
list is empty. This removes segments that cannot forward traffic in that time
period.

\subsection{Global maximum-load pruning}

If a finite incumbent upper bound $U$ is available, selecting segment $(i,j)$
for demand $d$ at time $t$ alone induces at least
\begin{equation}
\ell_{d,t,i,j}:=
\max_{a:r_{i,j,a,t}>0}
\frac{r_{i,j,a,t}\nu_{d,t}}{c_a}.
\tag{7}
\end{equation}
If
\[
\ell_{d,t,i,j}>U+\varepsilon_{\mathrm{lex}},
\]
the corresponding $x$ variable is not created because it cannot occur in a
solution whose global maximum utilization is at most $U$.

The pruning cutoff must always be a valid bound on the \emph{global} maximum
$L_1$, not the smaller maximum of the currently unfixed residual loads. A
segment may contribute more than a later residual boundary to an already fixed
earlier bottleneck. After the first full iteration, the optimal value $L_1$
is used to fix additional existing segment variables to zero without rebuilding
the persistent model.

\subsection{Load bounds during the solve}

The initial full model can use $U+\varepsilon_{\mathrm{lex}}$ as an upper bound
on its load variables and also adds $z\leq U$ when $U$ is available. During a
persistent bottleneck test, every still-free load is temporarily bounded by the
current simple-max boundary plus $\varepsilon_{\mathrm{lex}}$. Proven bottleneck
loads retain a narrow interval around their fixed value.

When the residual formulation is activated, the still-free loads receive the
explicit upper bound $R_1+\varepsilon_{\mathrm{lex}}$, while fixed bottlenecks
retain their own narrow bounds. The global $L_1$ cutoff continues to be used
for segment pruning.

\section{Lexicographic objective}

Let the $m$ utilization values be sorted in non-increasing order:
\[
L_1\geq L_2\geq\cdots\geq L_m,
\qquad
S_k:=\sum_{j=1}^{k}L_j.
\]
The target is lexicographic minimization of $(L_1,L_2,\ldots,L_m)$. The
implementation solves this sequentially, retaining the cumulative values
$S_k$. The ordered component can be recovered from
\[
L_1=S_1,
\qquad
L_k=S_k-S_{k-1}\quad(k\geq2).
\]

\subsection{First iteration: global min--max}

The first iteration minimizes the maximum load:
\begin{align}
S_1:=\min\quad &z \\
\text{s.t.}\quad
&z\geq\lambda_{a,t}, &&\forall(t,a)\in\Acal, \tag{8a}\\
&z\leq U, &&\text{if a finite warm-start bound is available}, \tag{8b}\\
&\text{constraints (1)--(6).} \nonumber
\end{align}
The cutoff in (8b) is optional and is only an incumbent bound; correctness does
not depend on the common-routing warm start succeeding.

\subsection{Persistent bottleneck peeling}

After each simple-max iteration, the implementation selects a candidate from
the \emph{physical} utilizations recomputed directly from the $x$ values. The
candidate is the largest still-unfixed load, rather than necessarily the largest
reported load variable in a model that might contain slack.

Suppose the current simple-max boundary is $L_k$ and the previously proved
bottlenecks are stored in a set $F$. The bottleneck test is performed on the
same native solver session and does not reconstruct the base model:

\begin{enumerate}[nosep]
    \item Every still-free load is temporarily bounded by
    $L_k+\varepsilon_{\mathrm{lex}}$. Already fixed loads keep their narrow
    fixed-value bounds.
    \item The objective coefficient of the persistent maximum variable $z$ is
    changed from $1$ to $0$, and the coefficient of the candidate load
    $\lambda_{\hat a,\hat t}$ is changed to $1$. A temporary objective link
    variable is not needed.
    \item The candidate load is minimized. A backend callback may stop as soon
    as it finds a feasible candidate strictly below
    $L_k-\tau$. Such an incumbent already disproves unavoidability, so a full
    optimality proof is unnecessary in that case.
    \item If the helper solve is optimal and its minimum is at least
    $L_k-\tau$, the candidate is declared individually unavoidable. If it can
    be reduced, or if the helper solve is interrupted or otherwise not proven
    optimal, no fixing is performed.
\end{enumerate}

The candidate is checked in three ways: the temporary objective value must agree
with the candidate $\lambda$ value, the candidate $\lambda$ must agree with the
physical utilization recomputed from $x$, and the helper optimum must not exceed
the feasible current boundary by more than the diagnostic tolerance.

If a candidate is proved unavoidable, the implementation:
\begin{enumerate}[nosep]
    \item fixes its load variable to a narrow interval around $L_k$;
    \item restores the objective to minimization of $z$; and
    \item removes the candidate's $\lambda$ coefficient from its persistent
    max-load row. The row $z-\lambda_{\hat a,\hat t}\geq0$ thereby becomes
    $z\geq0$, so the fixed load is peeled out of the next max objective while
    the row and column numbering remain unchanged.
\end{enumerate}

The same native model is then re-optimized for the next level. Consequently,
peeled iterations reuse the original rows, columns, presolve information, and
solver state as far as the backend permits. No new max variable is added for
each peeled level.

\subsection{Stopping peeling and switching to the residual formulation}

Let $m_f$ be the number of bottleneck loads that were proved individually
unavoidable when peeling stops. The simple-max iteration that detects the failed
proof has already optimized the first residual level
\[
R_1=L_{m_f+1}.
\]
The corresponding prefix objective is
\[
S_{m_f}=\sum_{j=1}^{m_f}L_j,
\]
with $S_0=0$. The next iteration is therefore not a global ordered-$k$-sum
model; it is a residual ordered-sum model over only the still-unfixed loads.

\subsection{Residual ordered-$k$-sum formulation}

Let $\Fcal$ denote the positions of the still-unfixed load variables, and let
$q=k-m_f$ be the residual order for global iteration $k$. The residual level
$R_1$ is preserved directly by the explicit bounds
\[
\lambda_p\leq R_1+\varepsilon_{\mathrm{lex}},
\qquad p\in\Fcal.
\]
For $q\geq2$, the model introduces $\alpha_i$ and $\beta_p^i$ only for
$i=2,\ldots,q$ and $p\in\Fcal$. The current residual objective is
\begin{align}
R_q:=\min\quad
&q\alpha_q+\sum_{p\in\Fcal}\beta_p^q \tag{9}\\
\text{s.t.}\quad
&\beta_p^i\geq\lambda_p-\alpha_i,
&&p\in\Fcal,\ i=2,\ldots,q, \tag{10}\\
&\beta_p^i\geq0,
&&p\in\Fcal,\ i=2,\ldots,q, \tag{11}\\
&i\alpha_i+\sum_{p\in\Fcal}\beta_p^i
\leq R_i+\varepsilon_{\mathrm{lex}},
&&i=2,\ldots,q-1, \tag{12}\\
&\alpha_i\in\R,
&&i=2,\ldots,q. \tag{13}
\end{align}
Fixed bottleneck loads are excluded from the $\alpha/\beta$ representation. Only
the current residual level $q$ receives objective coefficients; earlier residual
levels are preserved by (12). The global cumulative objective and current ordered
component are reconstructed as
\[
S_k=S_{m_f}+R_q,
\qquad
L_k=S_k-S_{k-1}.
\]
This residual formulation avoids repeatedly introducing auxiliary variables for
already fixed bottleneck loads and is smaller than a global ordered-$k$-sum model.

\section{Initial common-routing warm start}

Before the main lexicographic sequence, the implementation optionally solves a
smaller auxiliary model with one routing decision
\[
\bar{x}_{i,j}^{d}
\]
per demand and segment, shared by all time periods. A segment is eligible only
if it is reachable in every time period. Loads remain time-dependent and are
computed separately for every $(t,a)\in\Acal$ using the same exact load equations
as the full model.

The auxiliary model uses flow conservation, incoming/outgoing degree limits,
and the maximum-segment constraint once per demand. It minimizes its maximum
load. If a feasible solution is found, its routing is copied to every time
period as a MIP start for the first full model, and its measured maximum load
provides the valid bound $U$. If the auxiliary model fails, the full model is
solved without this warm start and without the associated cutoff.

Because the route is common to all time periods, it causes zero routing changes
in the reconfiguration constraints of the full model.

\section{Warm starts and persistent model reuse}

The first full model can receive the common-routing solution as a partial MIP
start. While bottleneck peeling is active, subsequent iterations are solved by
re-optimizing the same native solver session. The code changes column bounds,
objective coefficients, and one matrix coefficient per peeled bottleneck instead
of rebuilding the base model.

Once peeling has stopped, the residual models are rebuilt. Routing, load, and
reconfiguration values from the previous feasible solution are copied by
semantic position, so changes in column numbering caused by preprocessing do not
invalidate the warm start. For new residual auxiliary variables, the previous
free-load vector is sorted and a threshold is chosen from the corresponding
order statistic. The initialization is
\[
\beta_p^i=\max\{0,\lambda_p-\alpha_i\}.
\]
The solver remains free to improve this initialization.

\section{Model size}

Let $X$ and $Y$ be the sets of actually created segment and reconfiguration
variables, respectively, and let $m=|\Acal|$.

The initial full model contains $|X|$ binary variables, $m$ utilization
variables, $|Y|$ reconfiguration variables, and one maximum-load variable $z$.
During persistent peeling, no new columns are added. A proved bottleneck is
represented by a narrow load-variable bound and by removing one coefficient from
its existing max row.

If peeling stops after $m_f$ bottlenecks have been fixed, let $n_f$ be the number
of remaining free load positions and let $q$ be the residual order. A residual
model of order $q\geq2$ introduces
\[
(q-1)(n_f+1)
\]
additional continuous variables: one $\alpha_i$ and $n_f$ excess variables for
each residual level $i=2,\ldots,q$. No additional binary variables are introduced
by the lexicographic extensions.

Without sparse preprocessing, the segment-variable count is bounded by
$|D||T||V|(|V|-1)$ and the reconfiguration-variable count by
$|D||\Tstar||V|(|V|-1)$.

\section{Global time limit, interruption, and incumbent selection}

The command-line time limit is one global wall-clock budget for the complete
solve process. It includes the common-routing warm-start attempt, model
construction, all persistent re-optimizations, all bottleneck tests, and all
rebuilt residual models. Before each backend call, the remaining time is passed
to the solver.

If a time limit or a single Ctrl+C interruption prevents an iteration from
finishing, a feasible incumbent returned by that iteration is retained whenever
one exists. If no incumbent is returned, the best previously retained feasible
solution is used. The process compares feasible candidates using their actual
load vectors: both vectors are sorted in non-increasing order and compared
lexicographically up to the numerical tolerance. This allows a feasible
incumbent from an interrupted iteration to replace the last completed
lexicographic solution when it is better according to the intended load-vector
ordering.

A feasible common-routing solution can serve as a fallback if the main model has
not yet produced a full-model incumbent.

\section{Solver-independent backend}

The formulation is assembled in a generic sparse linear-model representation
containing variable names, objective coefficients, bounds, integrality types,
and sparse ranged rows. The representation is translated at runtime to either
backend:
\begin{itemize}[nosep]
    \item \textbf{HiGHS}: translated through the HiGHS C++ API;
    \item \textbf{Gurobi}: translated through the Gurobi C API.
\end{itemize}

The persistent backend interface supports changing column bounds, objective
coefficients, row bounds, and individual matrix coefficients. Both backends also
support the interrupt callback. The bottleneck helper uses an optional callback
criterion that stops minimization as soon as a feasible incumbent below the
rejection threshold is found.

The solver is selected at runtime with
\texttt{--solver highs} or \texttt{--solver gurobi}; \texttt{grb} is accepted as
an alias for Gurobi.

\section{LP model export}

With \texttt{--write-models 1}, the implementation exports the models passed to
the selected backend. Filenames include the solver and iteration information.
The common-routing model, persistent lexicographic models, bottleneck-test
re-optimizations, and residual ordered-sum models can be exported when the
corresponding stage is executed. The exported files are useful for inspecting
the exact bounds, coefficients, and variable names seen by the backend.

\section{Solution extraction and output}

For each demand and time period, selected binary segment variables define a
successor relation. Starting at the source, the extractor follows successors
until the target is reached and records the intermediate nodes as waypoints.
It rejects branching, cycles, missing successors, and paths exceeding the
maximum number of segments.

The resulting paths are written to the requested
\texttt{*-srpaths.json} file. Writing is also performed after a global timeout
or graceful Ctrl+C whenever a feasible incumbent is available.

\section{Program input and command-line interface}

For an instance named \texttt{setA-01}, the program expects
\begin{lstlisting}
instances/setA-01-net.json
instances/setA-01-tm.json
instances/setA-01-scenario.json
\end{lstlisting}
The split coefficients are computed from these files before model construction.

A typical invocation is
\begin{lstlisting}
./build/solver.exe \\
  --instance setA-01 \\
  --output setA-01-srpaths.json \\
  --iterations 10 \\
  --time 590 \\
  --solver highs \\
  --write-models 0 \\
  --model-prefix debug-setA-01
\end{lstlisting}

The relevant options are:
\begin{longtable}{>{\ttfamily}p{0.25\textwidth}p{0.63\textwidth}}
\toprule
\normalfont\textbf{Option} & \textbf{Meaning}\\
\midrule
--instance NAME & Instance basename.\\
--output FILE & Output JSON file containing the segment paths.\\
--iterations N & Maximum number of lexicographic iterations.\\
--time SECONDS & Global time budget for the complete solution process.\\
--solver highs|gurobi & Select the optimization backend at runtime.\\
--write-models 0|1 & Disable or enable LP export.\\
--model-prefix PREFIX & Prefix used for exported model files.\\
\bottomrule
\end{longtable}

\section{Windows development setup}

The current development setup uses \textbf{MSYS2 UCRT64}, CMake, C++20,
\texttt{networktools}, HiGHS installed through vcpkg for the MinGW triplet, and
optionally Gurobi. The project must be configured from the UCRT64 shell so that
an UCRT64 compiler is used.

For the HiGHS configuration used during development, the following is an
example command sequence:
\begin{lstlisting}
rm -rf build
cmake -B build \\
  -G "MinGW Makefiles" \\
  -DCMAKE_BUILD_TYPE=Release \\
  -DCMAKE_TOOLCHAIN_FILE=C:/vcpkg/scripts/buildsystems/vcpkg.cmake \\
  -DVCPKG_TARGET_TRIPLET=x64-mingw-dynamic
cmake --build build -j
\end{lstlisting}

The optional Gurobi backend uses the Gurobi C API. Set \texttt{GUROBI\_HOME} to
the installed \texttt{win64} directory and ensure that its \texttt{bin}
directory is on the MSYS2 \texttt{PATH}, for example:
\begin{lstlisting}
export GUROBI_HOME="/c/gurobi1302/win64"
export PATH="$GUROBI_HOME/bin:$PATH"
\end{lstlisting}

\section{Compact interpretation}

The implementation minimizes the sorted arc--time utilization vector in a
sequence of MILPs. The first model minimizes the global maximum load. While
possible, bottleneck peeling proves individual loads unavoidable by minimizing
one physical candidate under the current boundary, then fixes proven loads and
reuses the same native model for the next min--max level. If a proof fails, the
already optimized first residual maximum is preserved directly and later levels
are solved with a residual ordered-$k$-sum formulation that excludes fixed
loads. All load equations are exact, and physical-load diagnostics are used to
check the consistency of the solver model and the selected routing.

\end{document}
